\documentclass[11pt,a4paper]{article}
\usepackage[T1]{fontenc}
\usepackage{lmodern,amsmath,amssymb,booktabs,microtype}
\usepackage[margin=24mm]{geometry}
\usepackage[colorlinks=true,urlcolor=blue,linkcolor=blue]{hyperref}
\newcommand{\EE}{\mathbb E}
\newcommand{\tr}{\operatorname{tr}}
\newcommand{\W}[1]{\mathsf W_{#1}}
\newcommand{\Jh}{\mathcal J_h}
\newcommand{\Rem}{\operatorname{Rem}}
\title{Summing the charge sectors of the memory fluctuations\\From central windows to the infinite-temperature state}
\author{Alan Day \and Chengheng\thanks{AI research assistant (OpenAI).}}
\date{Local research note --- 6 November 2026}
\begin{document}\maketitle
\begin{abstract}
We extend the project's conserved-subtracted memory coefficients
to all charge sectors and control their charge tails. A uniform
Hamiltonian variance bound and a positive regression denominator
remove endpoint ambiguities. The central estimates extend to a
logarithmic window, while a coarse polynomial bound suffices
outside it. The full infinite-temperature spectral matrix equals
the central one. The rescaled charge is Gaussian and factors
from this two-energy matrix. This also supplies a precise
energy-diagonal subtraction for the subsequent finite-memory problem.
\end{abstract}
\section{The full state and the statement}
Use the same overlapping complex Gaussian $p$-body Hamiltonian,
standard CAR, $\hbar=1$, string-time units, and common coupling variance
\begin{equation}
 w_N=4^p\frac{J_s^2}{\lambda_N}\left(\frac{p!}{N^p}\right)^2,
 \qquad \lambda_N=p^2/N\to\lambda>0.
\end{equation}
$N$ is even, $\lambda_N\in[\lambda_-,\lambda_+]\subset(0,\infty)$,
and $p\le N/4$ eventually. The annealed state at $\beta=\mu=0$ is
\begin{equation}
 \EE\,2^{-N}\operatorname{Tr}X
 =\sum_{m=0}^N b_N(m)\EE\tr_mX_m,\qquad
 b_N(m)=2^{-N}\binom Nm,\quad x_m=(m-N/2)/p.
\end{equation}
The charge decomposition is standard in complex SYK~\cite{BNR};
the present question is uniformity for the \emph{subtracted}
memory fields, whose raw coefficients can grow with $p$.

Let $\mathbf V_m=(Y_m,L_m,U_{2,m})$ be the actual first three
collective memory coefficients with the conserved subtractions
of the central-raw note~\cite{central}. Their all-sector definition
and estimates are given below. If $d\boldsymbol\nu_*$ denotes
the central positive two-energy matrix established there, then
\begin{equation}
 \boxed{\quad d\boldsymbol\nu_N^{\rm full}
   =\sum_m b_N(m)d\boldsymbol\nu_{N,m}
   \ \Longrightarrow\ d\boldsymbol\nu_* .\quad}\label{full}
\end{equation}
Convergence is for bounded continuous energy tests, in particular
fixed real-time correlators. A stronger joint statement is
\begin{equation}
 d\boldsymbol\nu_N(x,E,E')\Longrightarrow
 \sqrt{\frac{2\lambda}{\pi}}e^{-2\lambda x^2}\,dx\,
 d\boldsymbol\nu_*(E,E').\label{joint}
\end{equation}
This is factorization of the indicated spectral matrix, not a
claim of independence of the complete finite-$N$ operator algebras.

\newpage
\section{Endpoint sectors and the variance bound}
Write $h=H_m$, $C_m=\EE\tr_mh^2$. The exact variance is
\begin{equation}
 C_m=w_N\binom mp\binom{N-m+p}p.
\end{equation}
For $m<p$, $H_m=H_{m-1}=0$ and the memory vanishes; all three
fields and their counterterms are defined to be zero. The
vacuum memory is likewise zero. For every $p\le m\le N$,
$C_m>0$. Term by term,
\begin{align}
 \frac{C_m}{J_s^2/\lambda_N}
 &=\prod_{j=0}^{p-1}\frac{4(m-j)(N-m+p-j)}{N^2}\notag\\
 &\le\prod_{j=0}^{p-1}\left(1+\frac{p-2j}{N}\right)^2
 \le e^{2p/N}.\label{maxC}
\end{align}
Indeed the difference of the two quadratic numerators is
$(2m-N-p)^2$, and $\sum_{j=0}^{p-1}(p-2j)=p$.
The same maximum is attained in the continuous charge variable
at $m=(N+p)/2$, up to integer rounding. Denote the right-hand
variance bound by $M_N$; it is bounded uniformly in the limit.
Also, exactly on active sectors,
\begin{equation}
 \frac{C_{m-1}}{C_m}
 =\frac{(m-p)(N-m+p+1)}{m(N-m+1)}\le p+1.\label{adj}
\end{equation}

Set $A_m=\sum_i[H,c_i]^\dagger[H,c_i]$, $B_m=h^2-C_mI$,
and $O_m=\W2 A_m$. The covariance eigenvalues are
$\kappa_a=w_N\binom{m-a}p\binom{N-m+p-a}p$ on the supported
rank subspaces. They lie in $[0,C_m]$. With their normalized
weights $\omega_a$, Wick contraction gives
\begin{equation}
 \|B_m\|_2^2=C_m^2(1+\chi_m),\qquad
 \chi_m=\sum_a\omega_a\kappa_a/C_m\in[0,1].\label{denom}
\end{equation}
Therefore $c_m=\langle O_m,B_m\rangle/\|B_m\|_2^2$ is well
defined even at $m=p,N$. No continuation through a missing
rank denominator is needed: use these invariant inner products.

We use normalized annealed Schatten norms below. Gaussian moments
obey $\EE\tr_m h^{2k}\le(2k-1)!!C_m^k$. One direct proof writes
$h=\sum g_iX_i$, $\sum X_i^2=C_mI$ and applies Gaussian
integration by parts. The map $\Phi(X)=\sum X_iXX_i$ divided
by $C_m$ is unital and trace preserving, hence contractive on
Schatten norms. H\"older bounds each of the $2k-1$ resulting
terms by $C_m\EE\tr_m|h|^{2k-2}$, completing the induction.

\newpage
\section{All-sector counterterms and a coarse global bound}
For $m\ge p$ use the central-raw formulas without rank cutoffs:
$a_m=\EE\tr_m A_m$, $K_A=a_mI+c_mB_m$, and
\begin{align}
 P_m&=(ac_m+b)h^3+[\bar t_m-(ac_m+b)\bar z_m]h,\notag\\
 \bar t_m&=\langle h,T_{3,m}\rangle/C_m,
 &\bar z_m&=\langle h,h^3\rangle/C_m.
\end{align}
Here $a,b$ are the fixed bounded functions of $q_N,\lambda_N$
in Ref.~\cite{central}. Define $r=p/m$ and
\begin{align}
 C_{0,m}&=K_A-mr^2h^2,\notag\\
 C_{1,m}&=P_m+2rhK_A-mr^3h^3,\notag\\
 D_{2,m}&=2rhP_m+3r^2h^2K_A-K_A^2/m-mr^4h^4,\notag\\
 C_{2,m}&=D_{2,m}+e_{\mathcal Q,m}I+\delta h^4
                 +g_Jh^2K_A+b_RhP_m\notag\\
 &\quad+C_m[p\ell_Q+k_Q+2\lambda_Nx_m\ell_Q
                  +o_Qc_m+8\lambda_Nq_N^2c]B_m.\label{C2}
\end{align}
All dimensionless coefficients and the raw polynomial
$\mathcal Q$ are exactly those defined in Ref.~\cite{central};
$e_{\mathcal Q,m}=\EE\tr_m\mathcal Q_m$ is computed before any
limit. These formulas are now a definition outside the central
window, with no assertion of a sectorwise asymptotic expansion there.
Every term commutes with $h$. For
$\widehat M_m(s)=G_m+isZ_{1,m}-s^2Z_{2,m}/2+O(s^3)$ put
\begin{equation}
 Y_m=mG_m-C_{0,m},\quad L_m=mZ_{1,m}-C_{1,m},\quad
 U_{2,m}=mZ_{2,m}-C_{2,m}.
\end{equation}

The normalized map $\Rem/m$, where $\Rem X=\sum c_i^\dagger Xc_i$,
is unital and preserves normalized traces from $m-1$ to $m$.
It is Schatten contractive. Thus the exact identity
$A_m=\Rem H_{m-1}^2+(2p-m)h^2$ gives
$\|A_m\|_{2k}\le K_kNM_N$. Equations \eqref{adj}--\eqref{denom}
also give $|c_m|\le KN^2$ and $\|K_A\|_{2k}\le K_kN^2M_N$.
The explicit three-contraction formula for $\bar t_m$ gives
$|\bar t_m|\le KN^2C_m$, while $|\bar z_m|\le3C_m$.
To see the first bound without small-variance division, each
binomial covariance ratio in $t_a/C_m$ is at most $p+1$ and
its reduced coefficient $\alpha_a$ lies in $[0,p]$.
Consequently $\|P_m\|_{2k}\le K_kN^2M_N^{3/2}$.

All raw quartic words in $\mathcal Q$ are bounded by
$K_kNM_N^2$, using $\Rem$, Gaussian moments and H\"older.
The exact renewal formula for $mZ_2$ has only these words,
products $hT_3,h^2A,A^2/m,h^4$, and coefficients bounded by
powers of $N$ with $r\le1$. In \eqref{C2} the largest crude
power is $K_A^2/m$. Therefore, with constants independent of $m$,
\begin{equation}
 \|Y_m\|_{2k}\le K_kN^2M_N,\quad
 \|L_m\|_{2k}\le K_kN^2M_N^{3/2},\quad
 \|U_{2,m}\|_{2k}\le K_kN^4M_N^2.\label{coarse}
\end{equation}
These deliberately coarse bounds are used only on exponentially
rare charge sectors, never to estimate the central fluctuations.

\newpage
\section{Auditing a logarithmically growing central window}
Fix any $A_0>0$ and set $R_N=\sqrt{A_0\log p}$.
The existing central calculation extends uniformly to
$|x_m|\le R_N$. The needed consequence is
\begin{equation}
 \sup_{|x_m|\le R_N}
 (\|Y_m\|_2^2+\|L_m\|_2^2+\|U_{2,m}\|_2^2)\le K_{A_0}.
 \label{innerbound}
\end{equation}
This is not inferred from compact-window uniformity alone.
The remainder audit is as follows.

\paragraph{Rational filters and the sixth moment.}
All scaled denominators of the exact filters stay separated from
zero, since $x_m/p=O(\sqrt{\log p}/p)$ and $a,b\le p,r\le2p$.
The finite sixth-moment recursion has rational coefficients in
$\epsilon=N^{-1/2},A,B,z=x_m\sqrt{\lambda_N}$.
The 21 negative-power cancellations are polynomial identities
in $z$, and its leading six moments are independent of $z$.
After cancellation, fixed-order Taylor remainders are bounded by
$K\epsilon(1+|z|)^D$ for a fixed finite $D$: each denominator
is a finite product of factors $1\pm2z\epsilon+O(\epsilon^2)$.
This gives a uniform sixth-moment bound on the logarithmic window.

\paragraph{The coherent force and input tails.}
In the exact coherent norm series, truncate input deficits
$u,v$ and the Racah index $j$ at $\lceil\log p\rceil$.
The input and summand ratios have fixed factorial envelopes
$K^u/u!$ and $K^j/j!$ throughout this window of charges.
The discarded terms, even multiplied by the old $p^4$ crude bound,
vanish faster than every power of $p$.
On the retained set the logarithmic expansion has a remainder
$p^{-3}$ times a fixed polynomial in $u,v,j,x_m$.
The exact $p^2,p$ cancellations and the $x_m$-independence of
the finite term therefore leave $p^{-1}(\log p)^D$.
This proves the central weighted $A$ norm bound here as well.

\paragraph{Strong substitutions.}
The $A$-filter error uses the sixth moment and an output cutoff
$S\le p^{3/4}$ as before. The $H^2$ filters subtract their exact
$p\ell$ and $2\lambda_Nx_m\ell$ terms first. On $S\le p^b$,
their Taylor remainder is $p^{-1}$ times a fixed polynomial;
choose fixed $b>0$ small enough for this polynomial degree.
On the complement, the linear filter bound and sixth moment
give an $L^2$ tail tending to zero. All extra charge factors are
powers of $\log p$. The cubic rank slope, centered $B_S,B_Z$,
and the $F$ proxy have the same factorial-weighted estimates.
The general force identity has errors proportional to $x_m/p$.
It replaces $\mathcal S_N$ by $(O_-+O_+)/2$ without a new
charge transport argument. Finite Gram projections are stable
because their limiting Gram is strictly positive at fixed $\lambda$.

\paragraph{Marked energy factors.}
The balanced-occupation trace proof gives error at most
$p^{-1}e^{K R_N}(1+R_N)^D$ per fixed marked graph after averaging
over overlaps; the same exponential overlap moments bound its tail.
This is $p^{-1+o(1)}$. It bounds all fixed energy insertions of
the marked representatives. The invariant Gaussian moment bound
controls multiplication of the linear filter errors by $h$.
Thus the restored raw expansion has a uniformly bounded leading
part and a vanishing strong remainder, proving \eqref{innerbound}.

\newpage
\section{Charge tails, full spectra, and joint factorization}
For $m\sim\operatorname{Bin}(N,1/2)$ the elementary Chernoff bound is
\begin{equation}
 \Pr(|x_m|>R)\le2e^{-2\lambda_N R^2}.
\end{equation}
For example it follows from
$\EE e^{t(m-N/2)}=\cosh(t/2)^N\le e^{Nt^2/8}$ and optimization
over $t$. Choose $A_0=20/\lambda_-$ in the logarithmic cutoff.
The squared version of \eqref{coarse} gives
\begin{equation}
 \sum_{|x_m|>R_N}b_N(m)\|\mathbf V_m\|_2^2
 \le K N^8e^{-2\lambda_-A_0\log p}
 \le K'p^{-24}\longrightarrow0.\label{outer}
\end{equation}
Here $\|\mathbf V\|_2^2$ is the sum of the three norm squares;
$M_N$ is absorbed into the constant. For a fixed $L<R_N$,
\eqref{innerbound} controls the intervening annulus:
\begin{equation}
 \sum_{L<|x_m|\le R_N}b_N(m)\|\mathbf V_m\|_2^2
 \le2K_{A_0}e^{-2\lambda_-L^2}.
\end{equation}
Together these establish the required uniform integrability:
\begin{equation}
 \lim_{L\to\infty}\limsup_{N\to\infty}
 \sum_{|x_m|>L}b_N(m)\|\mathbf V_m\|_2^2=0.\label{UI}
\end{equation}

For fields $V_i,V_j$, a bounded two-energy test has matrix element
bounded in modulus by its supremum norm times
$\|V_i\|_2\|V_j\|_2$. Thus \eqref{UI} controls both diagonal and
cross measures, including their total variation in the charge tails.
On $|x_m|\le L$ use the already proved compact-window spectral
limit, independent of $x_m$. First let $N\to\infty$ there,
then $L\to\infty$. This proves \eqref{full}; no sum over
energy-diagonal weights or oscillatory phases was discarded.

The charge characteristic function is exactly
\begin{equation}
 \EE e^{itx_m}=\cos\!\left(\frac{t}{2p}\right)^N,
 \qquad \log\EE e^{itx_m}
 =-\frac{t^2}{8\lambda_N}-\frac{t^4}{192\lambda_Np^2}
   +O(p^{-4}).
\end{equation}
It converges to the centered Gaussian with variance $1/(4\lambda)$.
For a bounded continuous charge test $f(x)$ and bounded continuous
energy test $g(E,E')$, the same compact-window argument gives
\begin{equation}
 \sum_m b_N(m)f(x_m)\int g\,d\boldsymbol\nu_{N,m}
 \longrightarrow
 \left[\int f(x)\sqrt{\frac{2\lambda}{\pi}}e^{-2\lambda x^2}dx\right]
 \int g\,d\boldsymbol\nu_* .
\end{equation}
Tightness and approximation on compact sets yield \eqref{joint}.
The charge fluctuation therefore does not add a new leading
dynamic spectral component once the correct sector-dependent
conserved polynomials have been removed. The ordered small-$\lambda$
results from the previous raw-spectrum note transfer to this full
state after the fixed-$\lambda$ double-scaled limit.

\newpage
\section{A definite finite-memory subtraction and executed evidence}
For each sample and charge let
\begin{equation}
 \mathsf P_h(X)=\sum_{\epsilon}P_\epsilon XP_\epsilon,
 \qquad h=\sum_\epsilon\epsilon P_\epsilon.
\end{equation}
This is the Hilbert--Schmidt orthogonal projection onto the
operators commuting with $h$, including degeneracies. It equals
the infinite symmetric time average of $e^{ith}Xe^{-ith}$ at
finite dimension. Define the actual finite-memory dynamic field
\begin{equation}
 \mathfrak F_{N,m}(s)=(1-\mathsf P_h)[m\widehat M_m(s)],\qquad
 \mathfrak F_{N,m}=0\quad(m<p).
\end{equation}
This subtracts the complete samplewise energy-diagonal part.
Its finite-$N$ Taylor coefficients through order two are exactly
$(1-\mathsf P_h)(Y,L,U_2)$, because all $C_{j,m}$ commute with $h$.
Moreover
\begin{equation}
 \|\mathsf P_hV\|_{2,\rm full}^2
 =\nu_{V,N}^{\rm full}\{E=E'\}\longrightarrow0,
 \quad V=Y,L,U_2.\label{pinch}
\end{equation}
The limiting positive diagonal measures in \eqref{full} are
absolutely continuous for $0<q<1$. Their mass on the closed
energy diagonal is zero, so weak convergence and the closed-set
inequality prove \eqref{pinch}. Consequently this exact subtraction
preserves the previously calculated limiting Taylor-coefficient
matrix. It specifies the next problem without guessing a
finite-$s$ polynomial background.

What remains is to bound and compute, for fixed real $s,u,T$,
$\EE\tr\mathfrak F_N(s)^\dagger e^{iTH}\mathfrak F_N(u)e^{-iTH}$.
The three known derivatives do not bound a finite-memory Taylor
remainder. The projection involves spectral information of each
sample and does not itself establish a local bulk reconstruction.

\paragraph{Checks.}
Exact binomial/product comparisons verify 6,766 sector variance
and adjacent-charge cases. Independent Gaussian-integer CAR
contractions cover every charge in the three systems
$(N,p)=(3,1)$, $(4,2)$ and $(5,2)$:
15 sectors, including inactive and fully filled ones. They verify
the denominator, $A$ mean, $B$--$O$ cross and regression residual.
Twelve new full-rank growing-window groups use
$x=\pm\lceil\sqrt{4\log p}\rceil$, $\lambda=0.25,0.5$,
$p=32,64,128$, at 80 and 120 digits. Four positive binomial-tail
sums check the outer cutoff. Six prior/current certificates and
their source hashes were verified; old scans were not rerun.
Finite-size differences and cancellation factors remain in the
JSON records. These numerical refinements are not interval bounds;
the logarithmic-window and summation proofs are the preceding sections.

\begin{thebibliography}{9}\small
\bibitem{BNR} M. Berkooz, V. Narovlansky and H. Raj,
\emph{Complex Sachdev--Ye--Kitaev model in the double scaling limit},
\href{https://arxiv.org/html/2006.13983v2}{arXiv:2006.13983v2},
Secs.~4.3 and 5. Charge projection and existing ensemble structure.
\bibitem{central} A. Day and Chengheng, local project notes:
\emph{The raw memory coefficient in a central charge window};
\emph{Coherent memory norms on central charge windows};
\emph{Central charge windows for the memory calculation};
\emph{Central force spectrum} (6 November 2026).
\end{thebibliography}
\end{document}
